\documentclass[letterpaper,11pt]{article}
\usepackage{graphicx}
\usepackage{geometry}
\linespread{1.2}                        
\setlength{\oddsidemargin}{0.0in}            
\setlength{\textwidth}{6.5in}    
\usepackage[utf8]{inputenc} % Asegúrate de usar utf8
\usepackage[T1]{fontenc} % Asegúrate de usar T1 font encoding      
\setlength{\topmargin}{-.6in}            
\setlength{\textheight}{9in}
\usepackage{latexsym} 
\usepackage{amssymb} 
\usepackage{amsmath}
\usepackage{amsxtra}
\usepackage[mathcal]{eucal} 
\usepackage{enumerate}
\usepackage{hyperref}
\usepackage{rotating}
\usepackage{caption}
\usepackage{lscape}
\usepackage{paralist} 
\usepackage{indentfirst}
\usepackage[dvipsnames,svgnames]{xcolor} 
\usepackage{setspace}
\usepackage{multicol} 
\usepackage{float}

\definecolor{codegreen}{rgb}{0,0.6,0}
\definecolor{codegray}{rgb}{0.5,0.5,0.5}
\definecolor{codepurple}{cmyk}{0.65, 1, 0, 0.35}
\definecolor{backcolour}{rgb}{0.95,0.95,0.92}

\usepackage{listings}
\usepackage{xcolor}

\lstset{
    language=R, % lenguaje de programación
    basicstyle=\small\ttfamily, % estilo básico
    commentstyle=\color{gray}, % estilo para los comentarios
    keywordstyle=\color{blue}, % estilo para las palabras clave
    stringstyle=\color{black}, % estilo para las cadenas de texto
    showstringspaces=false, % no mostrar espacios en cadenas de texto
    tabsize=4, % tamaño de la tabulación
    breaklines=true, % permitir el salto de línea automático
    postbreak=\mbox{\textcolor{red}{$\hookrightarrow$}\space}, % símbolo al final de una línea partida
    numbers=left, % mostrar números de línea
    numberstyle=\tiny\color{gray}, % estilo para los números de línea
    frame=single, % mostrar un borde alrededor del código
    backgroundcolor=\color{gray!10}, % color de fondo
}

\begin{document}

\begin{center}
    {\Large  Taller 3: Econometría }

Profesor:    Pedro Comte Hidalgo\\
Ayudante: Matias Balboa y Fabian Cisternas
\end{center}

\textbf{Instrucciones:} 

\begin{itemize}
    \item Usted dispone de hasta el día 5 de enero a las 21:00  para intentar los 100 puntos del taller.
    \item Cuide la presentación y redacción de sus respuestas.
    \item Puede utilizar su computador y los apuntes tanto de clase como de ayudantía.
    \item Debe enviar un archivo en formato PDF y un R script (extensión .R)
\end{itemize}

\section{Base de Datos htv (50 puntos)}

Para responder las preguntas siguientes, utilice la base de datos \textit{htv} de la librería \textbf{wooldridge}. Asegúrese de cargarla correctamente (sin nulos) y de especificar el programa de forma rigurosa para evitar cambios por aleatoriedad.

\begin{enumerate}

\item (10 puntos) Realice una matriz de correlación entre las variables de la base, identificando correlaciones altas, y eliminando, según criterio propio, una de las variables del par correlacionado.

\begin{lstlisting}
set.seed(1)
data(htv)
htv <- na.omit(htv)
matriz_corr <-cor(htv)
ggcorrplot(matriz_corr,type="full", lab=TRUE, lab_size = 3)
htv_filtrado <- subset(htv, select=-c(ne18,nc18,tuit17,lwage,expersq))
\end{lstlisting}


    \item (10 puntos) Realice un modelo de Regresión Lineal para las variables de esta base explicando el salario. Identifique y elimine de la base las variables que individualmente no son relevantes al 5\% de significancia. 
    
 \begin{lstlisting}
htv_filtrado <- subset(htv_filtrado, select=-c(ne,nc,motheduc,fatheduc,brkhme14,sibs,south18,west18,tuit18,ctuit,south,west))
modelo <- lm(wage ~.,data=htv_filtrado)
summary(modelo)
\end{lstlisting}   


    \item (30 puntos) Realice nuevamente un modelo de regresión lineal con la nueva base 
    \begin{enumerate}
        \item  (10 puntos) \textquestiondown Existe heterocedasticidad? Realice el test que más le acomode y concluya.
\begin{lstlisting}
bptest(modelo)  
#si existe heterocedasticidad
\end{lstlisting}
        \item (10 puntos) En caso de existir heterocedasticidad, encuentro los errores estandar robustos de White \textquestiondown Alguna variable que dejo de ser significante individualmente al 5\%?
\begin{lstlisting}
## Errores estándar robustos (White HC0)
coeftest(modelo, vcov = vcovHC(modelo, type = "HC0"))
## con estos se puede aplicar test t
#urban18 deja de ser relevante
\end{lstlisting}
        \item (10 puntos) Interprete los resultados del Modelo. \textquestiondown Tiene significancia Global? Considere que existe heterocedasticidad, utilizando la función \textit{waldtest} con solo un modelo y errores robustos. 
 \begin{lstlisting}       
#para hacer test F con heterocedasticidad hay que aplica lo siguiente
m_full <- lm(wage ~ abil +educ + exper + urban+urban18, data = htv)   # modelo sin restricciones

waldtest(m_full,
         vcov = vcovHC(m_full,  type = "HC0"),  # matriz robusta
         test = "F")
\end{lstlisting}        
    \end{enumerate} 
\end{enumerate}

\section{Base de Datos hprice2 (50 puntos)}

Para responder las preguntas siguientes, utilice la base de datos \textit{hprice2} de la librería \textbf{wooldridge}. Asegúrese de cargarla correctamente (sin nulos) y de especificar el programa de forma rigurosa para evitar cambios por aleatoriedad.

\begin{enumerate}

    \item (10 puntos) Realice una matriz de correlación entre las variables de la base, identificando correlaciones altas, y eliminando, según criterio propio, una de las variables del par correlacionado.

\begin{lstlisting}
data(hprice2)
hprice2 <- na.omit(hprice2)
matriz_corr <-cor(hprice2)
ggcorrplot(matriz_corr,type="full", lab=TRUE, lab_size = 3)

hprice2_filtrado <- subset(hprice2, select = -c(lprice,lnox,lproptax))
\end{lstlisting}
\item (10 puntos) Realice un modelo de Regresión Lineal para las variables de esta base explicando el salario. Identifique y elimine de la base las variables que individualmente no son relevantes al 10\% de significancia.  
\begin{lstlisting}
modelo <- lm(price ~., data=hprice2_filtrado)
summary(modelo)
\end{lstlisting}


   \item (30 puntos) Realice nuevamente un modelo de regresión lineal con la nueva base 
    \begin{enumerate}
        \item  (10 puntos) \textquestiondown Existe heterocedasticidad? Use test de White simplificado.
\begin{lstlisting}
bptest(modelo, ~ fitted(modelo) + I(fitted(modelo)^2)) 
#si existe heterocedasticidad
\end{lstlisting}
        \item (10 puntos) En caso de existir heterocedasticidad, encuentro los errores estandar robustos de White \textquestiondown Alguna variable que dejo de ser significante individualmente al 5\%?
\begin{lstlisting}
## Errores estándar robustos (White HC0)
coeftest(modelo, vcov = vcovHC(modelo, type = "HC0"))
## con estos se puede aplicar test t
#urban18 deja de ser relevante
\end{lstlisting}
        \item (10 puntos) Interprete los resultados del Modelo. \textquestiondown Tiene significancia Global? Considere que existe heterocedasticidad, utilizando la función \textit{waldtest} con solo un modelo y errores robustos. 
 \begin{lstlisting}       
#para hacer test F con heterocedasticidad hay que aplica lo siguiente
m_full <- lm(wage ~ abil +educ + exper + urban+urban18, data = htv)   # modelo sin restricciones

waldtest(m_full,
         vcov = vcovHC(m_full,  type = "HC0"),  # matriz robusta
         test = "F")
\end{lstlisting}        
    \end{enumerate} 

\end{enumerate}


\end{document}
